\ifdefined\gkver\else\def\gkver{student}\fi
\documentclass[\gkver]{gaokaozhenti}
\begin{document}

\chapter{1985年上海}

\begin{enumerate}
\item
%% number: 1
%% paperName: 1985年上海高考第 1 题 4 分
%% typeId: 3
%% type: 填空题
%% chapter: 电场
%% point: 库仑定律
%% method: 无
%% score: 4
%% degree: 850
%% duplicateId: 0
%% body:
在真空中有两个点电荷，电量分别是 $q_{1}$ 和 $q_{2}$，当它们相距为 $r$ 时，它们间相互作用力的大小已知是 $F$。如果它们的电量保持不变，而距离变为 $2r$，则它们之间的相互作用力的大小是\tkanswer{$\frac{1}{4}$}$F$。如果它们之间的距离 $r$ 保持不变，而电量分别变为 $2q_{1}$ 和 $3q_{2}$，则它们之间的相互作用力的大小是\tkanswer{6}$F$。

%% answer: 
\jdanswer{
$\frac{1}{4}$；6
}

%% memo:
\memoanswer{
本题原卷及材料中未提供详解，待补充。
}

\item
%% number: 2
%% paperName: 1985年上海高考第 2 题 4 分
%% typeId: 3
%% type: 填空题
%% chapter: 交流电
%% point: 交流电
%% method: 无
%% score: 4
%% degree: 850
%% duplicateId: 0
%% body:
照明用的交流电电压为 $220\UV$，这是指电压的\tkanswer{有效}值；当接入一个电阻为 $1100\UO$ 的电热丝时，通过电热丝的电流的最大值是\tkanswer{0.28}$\UA$。

%% answer: 
\jdanswer{
有效；0.28（如填 $0.2\sqrt{2}$ 也算对）
}

%% memo:
\memoanswer{
本题原卷及材料中未提供详解，待补充。
}

\item
%% number: 3
%% paperName: 1985年上海高考第 3 题 4 分
%% typeId: 3
%% type: 填空题
%% chapter: 机械振动
%% point: 振动图象
%% method: 无
%% score: 4
%% degree: 850
%% duplicateId: 0
%% body:
一个质点的简谐振动图象如图所示。从图中可以看出，振动的振幅是\tkanswer{0.10}$\Um$，频率是\tkanswer{2.5}$\UHz$。
\onepicture[align=right, w=6cm]{\includesvg{figs/03}}

%% answer: 
\jdanswer{
0.10；2.5
}

%% memo:
\memoanswer{
本题原卷及材料中未提供详解，待补充。
}

\item
%% number: 4
%% paperName: 1985年上海高考第 4 题 4 分
%% typeId: 3
%% type: 填空题
%% chapter: 原子和原子核
%% point: 玻尔模型
%% method: 无
%% score: 4
%% degree: 850
%% duplicateId: 0
%% body:
设氢原子基态的能量是 $E_{1}$，某一激发态的能量是 $E_{2}$，当氢原子从这一激发态跃迁到基态时，它辐射出的光子的频率 $\nu =$\tkanswer[1.6cm]{$\frac{E_{2}-E_{1}}{h}$}；在真空中的波长 $\lambda =$\tkanswer[1.8cm]{$\frac{hc}{E_{2}-E_{1}}$}。（光在真空中的速度记作 $c$，普朗克恒量记作 $h$。）

%% answer: 
\jdanswer{
$\frac{E_{2}-E_{1}}{h}$；$\frac{hc}{E_{2}-E_{1}}$（如填 $\frac{c}{\nu}$ 算对）
}

%% memo:
\memoanswer{
本题原卷及材料中未提供详解，待补充。
}

\item
%% number: 5
%% paperName: 1985年上海高考第 5 题 4 分
%% typeId: 3
%% type: 填空题
%% chapter: 原子和原子核
%% point: 半衰期
%% method: 无
%% score: 4
%% degree: 850
%% duplicateId: 0
%% body:
\begin{enumerate}
	\item
	平衡下列核反应方程式：$\ce{^{9}_{4}Be + ^{4}_{2}He -> ^{12}_{6}C}$ +\tkanswer[1.2cm]{$\ce{^{1}_{0}n}$}。
	\item
	铋 210 的半衰期是 $5\Ud$，$12\Ug$ 的铋 210 经过 $20\Ud$ 以后还剩下\tkanswer{0.75}$\Ug$。
\end{enumerate}

%% answer: 
\jdanswer{
\begin{enumerate}
	\item
	$\ce{^{1}_{0}n}$

	\item
	0.75

\end{enumerate}
}

%% memo:
\memoanswer{
本题原卷及材料中未提供详解，待补充。
}

\item
%% number: 6
%% paperName: 1985年上海高考第 6 题 3 分
%% typeId: 1
%% type: 单选题
%% chapter: 功和能
%% point: 功
%% method: 无
%% score: 3
%% degree: 850
%% duplicateId: 0
%% body:
水平地面上有一块重量是 $2\UN$ 的静止石块。一个小孩用 $10\UN$ 的力踢石块，使石块滑行了 $1\Um$ 的距离，则小孩对石块所做的功是\xzanswer{D}
\fourchoices[answer=D]
{$10\UJ$}
{$2\UJ$}
{$12\UJ$}
{条件不足，无法确定}

%% answer: D
%% memo:
\memoanswer{
本题原卷及材料中未提供详解，待补充。
}

\item
%% number: 7
%% paperName: 1985年上海高考第 7 题 3 分
%% typeId: 1
%% type: 单选题
%% chapter: 气体性质
%% point: 气体的状态参量
%% method: 无
%% score: 3
%% degree: 650
%% duplicateId: 0
%% body:
对于一定质量的理想气体，正确的是\xzanswer{B}
\fourchoices[answer=B]
{如果体积 $V$ 减小，气体分子在单位时间内作用于器壁单位面积的总冲量一定增大}
{如果压强 $p$ 增大，气体分子在单位时间内作用于器壁单位面积的总冲量一定增大}
{如果温度 $T$ 不变，气体分子在单位时间内作用于器壁单位面积的总冲量一定不变}
{如果密度不变，气体分子在单位时间内作用于器壁单位面积的总冲量一定不变}

%% answer: B
%% memo:
\memoanswer{
本题原卷及材料中未提供详解，待补充。
}

\item
%% number: 8
%% paperName: 1985年上海高考第 8 题 3 分
%% typeId: 1
%% type: 单选题
%% chapter: 电场
%% point: 电场线
%% method: 无
%% score: 3
%% degree: 650
%% duplicateId: 0
%% body:
有一电场的电场线分布如图所示，场中 A、B 两点的电场强度的大小和电势分别用 $E_{A}$、$E_{B}$ 和 $\varphi_{A}$、$\varphi_{B}$ 表示。则有\xzanswer{D}
\onepicture[align=right, w=4.5cm]{\includesvg{figs/08}}
\fourchoices[answer=D]
{$E_{A}$ 大于 $E_{B}$，$\varphi_{A}$ 高于 $\varphi_{B}$}
{$E_{A}$ 大于 $E_{B}$，$\varphi_{A}$ 低于 $\varphi_{B}$}
{$E_{A}$ 小于 $E_{B}$，$\varphi_{A}$ 高于 $\varphi_{B}$}
{$E_{A}$ 小于 $E_{B}$，$\varphi_{A}$ 低于 $\varphi_{B}$}

%% answer: D
%% memo:
\memoanswer{
本题原卷及材料中未提供详解，待补充。
}

\item
%% number: 9
%% paperName: 1985年上海高考第 9 题 3 分
%% typeId: 1
%% type: 单选题
%% chapter: 光学
%% point: 光电效应
%% method: 无
%% score: 3
%% degree: 650
%% duplicateId: 0
%% body:
用频率是 $\nu_{1}$ 的光照射某金属表面时，正好发生光电效应。如果改用强度相同、但频率是 $\nu_{2}$ 的另一种光（已知 $\nu_{2} > \nu_{1}$）照射这一金属表面，则\xzanswer{B}
\fourchoices[answer=B]
{光电子的最大初动能不变}
{光电子的最大初动能增大}
{光电子的最大初动能减小}
{不发生光电效应}

%% answer: B
%% memo:
\memoanswer{
本题原卷及材料中未提供详解，待补充。
}

\item
%% number: 10
%% paperName: 1985年上海高考第 10 题 3 分
%% typeId: 2
%% type: 多选题
%% chapter: 运动和力
%% point: 牛顿第三定律
%% method: 无
%% score: 3
%% degree: 850
%% duplicateId: 0
%% body:
粗糙的水平地面上有一木箱。现用一水平力拉着木箱匀速前进，则\xzanswer{BC}
\fourchoices[answer=BC]
{木箱所受的拉力和地面对木箱的摩擦力是一对作用力和反作用力}
{木箱对地面的压力和地面对木箱的支持力是一对作用力和反作用力}
{木箱所受的重力和地面对木箱的支持力是一对平衡力}
{木箱对地面的压力和地面对木箱的支持力是一对平衡力}

%% answer: BC
%% memo:
\memoanswer{
本题原卷及材料中未提供详解，待补充。
}

\item
%% number: 11
%% paperName: 1985年上海高考第 11 题 3 分
%% typeId: 2
%% type: 多选题
%% chapter: 磁场
%% point: 洛伦兹力
%% method: 无
%% score: 3
%% degree: 650
%% duplicateId: 0
%% body:
电子以初速 $v_{0}$ 垂直进入磁感应强度为 $B$ 的匀强磁场中，则\xzanswer{BD}
\fourchoices[answer=BD]
{磁场对电子的作用力始终不变}
{磁场对电子的作用力始终不做功}
{电子的动量始终不变}
{电子的动能始终不变}

%% answer: BD
%% memo:
\memoanswer{
本题原卷及材料中未提供详解，待补充。
}

\item
%% number: 12
%% paperName: 1985年上海高考第 12 题 3 分
%% typeId: 2
%% type: 多选题
%% chapter: 其他
%% point: 物理学史
%% method: 无
%% score: 3
%% degree: 650
%% duplicateId: 0
%% body:
下列叙述中，符合物理学史事实的有\xzanswer{ABC}
\fourchoices[answer=ABC]
{托马斯$\cdot$杨通过对光的干涉的研究，证实了光具有波动性}
{爱因斯坦为了解释光电效应的规律，提出了光子说}
{卢瑟福通过对 $\alpha$ 粒子散射的研究，提出了原子的核式结构学说}
{贝克勒耳通过对天然放射性的研究，发现了原子核是由质子和中子组成的}

%% answer: ABC
%% memo:
\memoanswer{
本题原卷及材料中未提供详解，待补充。
}

\item
%% number: 13
%% paperName: 1985年上海高考第 13 题 3 分
%% typeId: 2
%% type: 多选题
%% chapter: 电路
%% point: 闭合电路欧姆定律
%% method: 无
%% score: 3
%% degree: 450
%% duplicateId: 0
%% body:
如右图所示，一内部具有电路结构的盒子上有两个插孔。已知如把伏特表（内阻很大）正确接上插孔，读数是 $3\UV$；如把安培表（内阻可以忽略）正确接上插孔，读数是 $3\UA$；则下列画有内部电路结构的盒子中，哪几个是可能的？\xzanswer{AD}
\onepicture[align=right, w=4.5cm]{\includesvg{figs/13a}}
\fourchoices[answer=AD, ispicture=true, h=2.2cm]
{figs/13b.png}
{figs/13c.png}
{figs/13d.png}
{figs/13e.png}

%% answer: AD
%% memo:
\memoanswer{
本题原卷及材料中未提供详解，待补充。
}

\item
%% number: 14
%% paperName: 1985年上海高考第 14 题 4 分
%% typeId: 4
%% type: 实验题
%% chapter: 气体性质
%% point: 盖吕萨克定律
%% method: 无
%% score: 4
%% degree: 650
%% duplicateId: 0
%% body:
如图所示，在球形烧瓶上连一根水平玻璃管，管中装有一小段水银柱，用来把烧瓶中的气体和外界的大气隔开。先把烧瓶放进盛着冰水混合物的容器里，经过一段时间，再把烧瓶放进热水中，可以看到此时水平玻璃管中的水银柱将\tkanswer{向右}移动。在这个实验中，烧瓶中气体初、末状态参量的变化是：压强\tkanswer{不变}，温度\tkanswer{升高}，体积\tkanswer{增大}。
\onepicture[align=right, w=6cm]{\includesvg{figs/14}}

%% answer: 
\jdanswer{
向右；不变；升高；增大。
}

%% memo:
\memoanswer{
向右；不变；升高；增大。（各 1 分，共 4 分。）
}

\item
%% number: 15
%% paperName: 1985年上海高考第 15 题 6 分
%% typeId: 4
%% type: 实验题
%% chapter: 电路
%% point: 测电动势与内阻
%% method: 无
%% score: 6
%% degree: 650
%% duplicateId: 0
%% body:
在用伏特表和安培表测定电池的电动势和内电阻的实验中，所用的安培表量程有 $0\sim0.6\UA$ 和 $0\sim3\UA$ 两档。
\begin{enumerate}
	\item
	在连接好电路并按下电键以前，安培表应先接在\tkanswer{$0\sim3\UA$}档，滑动变阻器的滑动触片应先处在\tkanswer{电阻最大}位置上。
	\item
	如果根据实验数据画出的 $U-I$ 关系图线如图所示。则可求得电池的电动势是\tkanswer{1.5}$\UV$，内电阻是\tkanswer{0.75}$\UO$。理由是：\tkanswer{由 $U=E-Ir$ 可知：当 $I=0$ 时，$U=E$，即纵坐标 $U$ 轴上的截距等于 $E$，把图线延长，读出 $E=1.5\UV$；当 $U=0$ 时，$I=\frac{E}{r}$，即横坐标轴上的截距是 $\frac{E}{r}$，把图线延长，读出此电流是 $2.0\UA$，求得 $r=0.75\UO$}。
\end{enumerate}
\onepicture[align=right, w=6.5cm]{\includesvg{figs/15}}

%% answer: 
\jdanswer{
\begin{enumerate}
	\item
	$0\sim3\UA$；电阻最大

	\item
	$1.5\UV$；$0.75\UO$，理由：由 $U=E-Ir$ 可知：当 $I=0$ 时，$U=E$，即纵坐标 $U$ 轴上的截距等于 $E$，把图线延长，读出 $E=1.5\UV$；当 $U=0$ 时，$I=\frac{E}{r}$，即横坐标轴上的截距是 $\frac{E}{r}$，把图线延长，读出此电流是 $2.0\UA$，求得 $r=0.75\UO$。

\end{enumerate}
}

%% memo:
\memoanswer{
（1）$0\sim3\UA$（1 分。如填 $3\UA$ 也算对。）；电阻最大（1 分）。

（2）$1.5\UV$（1 分。可允许在 $1.4\sim1.6\UV$ 范围内。）；$0.75\UO$（1 分。可允许在 $0.70\sim0.80\UO$ 范围内。）理由：由 $U=E-Ir$ 可知：当 $I=0$ 时，$U=E$，即纵坐标 $U$ 轴上的截距等于 $E$，把图线延长，读出 $E=1.5\UV$；当 $U=0$ 时，$I=\frac{E}{r}$，即横坐标轴上的截距是 $\frac{E}{r}$，把图线延长，读出此电流是 $2.0\UA$，求得 $r=0.75\UO$。（2 分）
}

\item
%% number: 16
%% paperName: 1985年上海高考第 16 题 8 分
%% typeId: 4
%% type: 实验题
%% chapter: 功和能
%% point: 动能定理
%% method: 无
%% score: 8
%% degree: 650
%% duplicateId: 0
%% body:
给你一架天平和一只秒表，你如何用实验来估测用手急速竖直上抛小球时，手对小球所做的功？（要求：写出你准备测定的那些量，如何测量，并列出计算时所需的关系式。）

%% answer: 
\jdanswer{
测量小球的质量 $m$，用秒表测出小球从抛出到落回原处所经历的时间 $T$，由 $W=\frac{1}{8}mg^{2}T^{2}$ 估算手对小球所做的功。
}

%% memo:
\memoanswer{
测量小球的质量。（1 分）

先调节天平，然后称出小球的质量 $m$。（1 分，如未提到先调节天平的，可不扣分。）

测量小球从抛出到落回原处时所经历的时间。（1 分。如答测量小球从抛出到最高处的时间，不给分。）

用秒表测出这一时间 $T$。（1 分）

设小球被抛出时的初速度是 $v_{0}$，由竖直上抛运动的规律 $T=\frac{2v_{0}}{g}$，得 $v_{0}=\frac{1}{2}gT$。（2 分）

设手对小球所做的功是 $W$，由外力做功跟物体动能变化的关系，得

$W=\frac{1}{2}mv_{0}^{2}=\frac{1}{8}mg^{2}T^{2}$。（2 分。如只答到 $W=\frac{1}{2}mv_{0}^{2}$，给 1 分。）

注：用手上抛小球，在球离手前，先有一个使球加速上升的过程。这一过程的上升高度设为 $\Delta h$。则计算手对小球所做的功时，还应加上 $mg\Delta h$ 一项，在用手急速上抛小球时，$\Delta h$ 较小，又因是估测，可不计及。考生如能指出这一点或提及还需估算 $\Delta h$，则是很好的。
}

\item
%% number: 17
%% paperName: 1985年上海高考第 17 题 8 分
%% typeId: 6
%% type: 计算题
%% chapter: 运动和力
%% point: 牛顿定律的应用
%% method: 无
%% score: 8
%% degree: 650
%% duplicateId: 0
%% body:
一木块从高 $h=3.0\Um$、长 $l=5.0\Um$ 的固定斜面的顶端，由静止开始沿着斜面滑至底端。如果木块与斜面之间的滑动摩擦系数 $\mu=0.30$，求：
\begin{enumerate}
	\item
	木块运动的加速度；
	\item
	木块从斜面顶端滑至底端所需的时间。
\end{enumerate}
\onepicture[align=right, w=5cm]{\includesvg{figs/17}}

%% answer: 
\jdanswer{
\begin{enumerate}
	\item
	$a=3.6\Umsq$

	\item
	$t=1.7\Us$（或 $\frac{5}{3}\Us$）

\end{enumerate}
}

%% memo:
\memoanswer{
\onepicture[align=right, w=4cm]{\includesvg{figs/17_an}}

（1）解：如图所示，木块在斜面上运动时受重力 $G$（等于 $mg$）、支持力 $N$、摩擦力 $f$ 的作用。将重力 $G$ 沿跟斜面平行方向和垂直方向分解，由牛顿定律得：

$mg\sin\alpha-f=ma$　（1）（1 分）

$N-mg\cos\alpha=0$　（2）

其中 $\alpha$ 是斜面的倾角，由题中已知条件有：$\sin\alpha=\frac{h}{l}=0.60$，$\cos\alpha=0.80$。

由滑动摩擦力公式得：$f=\mu N=\mu mg\cos\alpha$　（3）（1 分）

代入得：$mg\sin\alpha-\mu mg\cos\alpha=ma$　（4）（1 分）

解得：$a=(\sin\alpha-\mu\cos\alpha)g=(0.60-0.30\times0.80)\times10\Umsq=3.6\Umsq$。（1 分）

（如直接写出（4）式也给 3 分。）

（2）解：木块沿斜面滑下时做匀加速运动，设所需时间为 $t$，由运动学公式 $s=\frac{1}{2}at^{2}$（1 分）

得 $t=\sqrt{\frac{2s}{a}}=\sqrt{\frac{2\times5.0}{3.6}}\Us$（1 分）

$=1.7\Us$（或 $\frac{5}{3}\Us$）（1 分）
}

\item
%% number: 18
%% paperName: 1985年上海高考第 18 题 8 分
%% typeId: 6
%% type: 计算题
%% chapter: 气体性质
%% point: 玻意耳定律
%% method: 无
%% score: 8
%% degree: 650
%% duplicateId: 0
%% body:
如图所示，把一根两端开口、粗细均匀的玻璃管竖直插入水银槽中，当玻璃管露出水银部分长 $27\Ucm$ 时，将上端封闭。然后把玻璃管缓慢地竖直压下 $8\Ucm$，问：
\begin{enumerate}
	\item
	这时管内水银面将如何变化？为什么？
	\item
	试求管内外水银面的高度差。（已知大气压强是 $75\UcmHg$）
\end{enumerate}
\onepicture[align=right, w=5cm]{\includesvg{figs/18}}

%% answer: 
\jdanswer{
\begin{enumerate}
	\item
	答：管中水银面将缓慢降低。把管缓慢地竖直压下时，如果水银面不动，管内空气体积将减小，而温度不变，从玻意耳$-$马略特定律可知，它的压强将增大，大于外界大气压强，所以管中水银面实际上将降低到另一平衡位置。

	\item
	$h=6\Ucm$

\end{enumerate}
}

%% memo:
\memoanswer{
（1）答：管中水银面将缓慢降低。

把管缓慢地竖直压下时，如果水银面不动，管内空气体积将减小，而温度不变，从玻意耳$-$马略特定律可知，它的压强将增大，大于外界大气压强，所以管中水银面实际上将降低到另一平衡位置。（3 分。只答管中水银面降低的，给 1 分。）

（2）解：设管内外水银面的高度差为 $h\Ucm$，由（1）可知，是管外水银面高于管内，又设管的截面积为 $S\Ucmq$，则：

管内空气的初状态：压强 $p_{1}=75\UcmHg$，体积 $V_{1}=27S\Ucmc$

管内空气的末状态：压强 $p_{2}=(75+h)\UcmHg$，体积 $V_{2}=(27-8+h)S\Ucmc$

由玻意耳$-$马略特定律 $p_{1}V_{1}=p_{2}V_{2}$（1 分）

代入得：$75\times27S=(75+h)(27-8+h)S$（3 分）

$h^{2}+94h-600=0$

$(h+100)(h-6)=0$

$h=6\Ucm$（1 分）

$h=-100\Ucm$（不合题意舍去）
}

\item
%% number: 19
%% paperName: 1985年上海高考第 19 题 10 分
%% typeId: 6
%% type: 计算题
%% chapter: 运动和力
%% point: 动量守恒定律
%% method: 无
%% score: 10
%% degree: 550
%% duplicateId: 0
%% body:
一质量 $m=20\Ukg$ 的小车停放在光滑固定轨道 ABCDE 的 A 点处，如图所示，轨道的 AB 段是水平的，CDE 段是一半径 $R=32\Um$ 的圆弧，圆弧的最高点 D 比 AB 处高出 $h=1.0\Um$。有一质量 $M=60\Ukg$ 的人以 $v_{0}=8.0\Ums$ 的水平速度跳上小车，并与小车一起沿轨道滑行。不计一切阻力，试计算：
\begin{enumerate}
	\item
	当人跳上小车后，小车的运动速度 $u$；
	\item
	当小车滑行到 D 点时小车对轨道的压力 $N$。
\end{enumerate}
\onepicture[align=right, w=8cm]{\includesvg{figs/19}}

%% answer: 
\jdanswer{
\begin{enumerate}
	\item
	$u=6.0\Ums$（方向同 $v_{0}$，向前）

	\item
	$N=760\UN$

\end{enumerate}
}

%% memo:
\memoanswer{
（1）由动量守恒定律：$Mv_{0}=(M+m)u$（1 分）

$u=\frac{Mv_{0}}{M+m}=\frac{60\times8.0}{60+20}\Ums$（1 分）

$=6.0\Ums$（方向同 $v_{0}$，向前）（1 分）

（2）设小车滑行到 D 点时的速度是 $v$，由机械能守恒定律，有：

$\frac{1}{2}(M+m)u^{2}=\frac{1}{2}(M+m)v^{2}+(M+m)gh$（1 分）

$v=\sqrt{u^{2}-2gh}=\sqrt{6^{2}-2\times10\times1.0}\Ums$（1 分）

$=4.0\Ums$（1 分）

小车在 D 点所受重力 $G$【等于 $(M+m)g$】和支持力 $N$ 的合力就是向心力，有：

$(M+m)g-N=\frac{(M+m)v^{2}}{R}$（1 分）

$N=(M+m)(g-\frac{v^{2}}{R})=(60+20)(10-\frac{16}{32})\UN$（1 分）

$=760\UN$（1 分）

小车对轨道的压力跟小车所受的支持力大小相等，方向相反，即大小也为 $760\UN$。（如只算小车质量，其余部分都正确而算出压力是 $190\UN$ 的，该最后 4 分中只给 2 分。）
}

\item
%% number: 20
%% paperName: 1985年上海高考第 20 题 12 分
%% typeId: 6
%% type: 计算题
%% chapter: 电磁感应
%% point: 电磁感应中的变加速
%% method: 无
%% score: 12
%% degree: 650
%% duplicateId: 0
%% body:
图中所示是一个水平放置的导体框架，宽度 $l=0.50\Um$，接有电阻 $R=0.20\UO$。设匀强磁场与框架平面垂直，磁感应强度 $B=0.40\UT$，方向如图所示。今有一条形导体 ab 跨放在框架上，并能无摩擦地沿框架滑动，框架和导体 ab 的电阻均不计。当 ab 以 $v_{0}=4.0\Ums$ 的速度向右匀速滑动时，
\begin{enumerate}
	\item
	求导体 ab 上的感应电动势的大小；
	\item
	求回路上感应电流的大小，并在图上标出电流的方向；
	\item
	要维持导体 ab 作匀速运动，必须要有外力 $F$ 作用在导体 ab 上，为什么？试求出 $F$ 的大小和方向；
	\item
	如果作用在导体 ab 上的外力突然减小为 $F'$，试简要讨论导体 ab 以后的运动情况。
\end{enumerate}
\onepicture[align=right, w=6cm]{\includesvg{figs/20}}

%% answer: 
\jdanswer{
\begin{enumerate}
	\item
	$E=0.80\UV$

	\item
	$I=4.0\UA$，电流方向是 aRba。

	\item
	答：导体 ab 在如图磁场中向右运动时，闭合电路中就有沿 aRba 的感生电流。此时，磁场对 ab 有作用力 $f$，按左手定则判定，$f$ 的方向是向左的。为了维持 ab 作匀速运动，所以必须加一个跟 $f$ 方向相反、大小相等的外力 $F$。

	\item
	当加在导体 ab 上的外力突然减小为 $F'$ 时，由于 $F'<F$，ab 将做减速运动，$v$ 减小，当 $v$ 减小时，由 $E=Blv$ 可知 $E$ 同时减小；由 $I=\frac{E}{R}$ 和 $f=BIl$ 可知 $I$ 和 $f$ 也都同时在减小；但只要 $F'<f$，则 $f$ 将继续减小，直到两力平衡，这时 ab 导体将以某一比 $v$ 小的速度重新做匀速运动。

\end{enumerate}
}

%% memo:
\memoanswer{
（1）解：由感生电动势大小公式 $E=Blv=0.40\times0.50\times4.0\UV$（1 分）

$=0.80\UV$（1 分）

（2）解：由欧姆定律：$I=\frac{E}{R}=\frac{0.80}{0.20}\UA=4.0\UA$（1 分）

应在图中标明电流方向是 aRba。（此处图略）（1 分）

（3）答：导体 ab 在如图磁场中向右运动时，闭合电路中就有沿 aRba 的感生电流。此时，磁场对 ab 有作用力 $f$，按左手定则判定，$f$ 的方向是向左的。为了维持 ab 作匀速运动，所以必须加一个跟 $f$ 方向相反、大小相等的外力 $F$。（2 分）

解：磁场对导体 ab 的作用力

$f=BIl$

$F=BIl=4.0\times0.50\times0.40\UN=0.80\UN$（1 分）

力的方向向右（1 分）

（4）答：当加在导体 ab 上的外力突然减小为 $F'$ 时，由于 $F'<F$，ab 将做减速运动，$v$ 减小，当 $v$ 减小时，由 $E=Blv$ 可知 $E$ 同时减小；由 $I=\frac{E}{R}$ 和 $f=BIl$ 可知 $I$ 和 $f$ 也都同时在减小；但只要 $F'<f$，则 $f$ 将继续减小，直到两力平衡，这时 ab 导体将以某一比 $v$ 小的速度重新做匀速运动。（4 分）
}

\end{enumerate}

\end{document}
